\documentclass[10pt,a4paper]{amsart}
\usepackage{amsmath,amssymb,amsthm,graphicx}
% Retained minimal preamble for the Stacks Project chapter more-algebra.tex (More on Algebra).
% The chapter file pinned-original-source.tex carries only \input{preamble} on its first line, so no
% preamble text lives inside the chapter itself.  The Stacks Project uses one single shared preamble for
% every chapter; the snapshot stages it identically as sources/stacks-chow/preamble.tex and
% sources/stacks-etale-cohomology/preamble.tex.  The block below is copied mechanically, line for line,
% from sources/stacks-chow/preamble.tex lines 14-18 (the Xy-pic setup) and 154-217 (the native theorem
% declarations together with the mathematical macros).  Nothing was added, removed or redefined; the
% wrapper class and the amsmath/amssymb/amsthm/graphicx base packages are supplied by the assembler.
% Chapter: more-algebra.tex -- More on Algebra.
% For commutative diagrams we use Xy-pic
\usepackage[all]{xy}

% We use 2cell for 2-commutative diagrams.
\xyoption{2cell}

% Theorem environments.
%
\theoremstyle{plain}
\newtheorem{theorem}[subsection]{Theorem}
\newtheorem{proposition}[subsection]{Proposition}
\newtheorem{lemma}[subsection]{Lemma}

\theoremstyle{definition}
\newtheorem{definition}[subsection]{Definition}
\newtheorem{example}[subsection]{Example}
\newtheorem{exercise}[subsection]{Exercise}
\newtheorem{situation}[subsection]{Situation}

\theoremstyle{remark}
\newtheorem{remark}[subsection]{Remark}
\newtheorem{remarks}[subsection]{Remarks}

\numberwithin{equation}{subsection}

% Macros
%
\def\lim{\mathop{\mathrm{lim}}\nolimits}
\def\colim{\mathop{\mathrm{colim}}\nolimits}
\def\Spec{\mathop{\mathrm{Spec}}}
\def\Hom{\mathop{\mathrm{Hom}}\nolimits}
\def\Ext{\mathop{\mathrm{Ext}}\nolimits}
\def\SheafHom{\mathop{\mathcal{H}\!\mathit{om}}\nolimits}
\def\SheafExt{\mathop{\mathcal{E}\!\mathit{xt}}\nolimits}
\def\Sch{\mathit{Sch}}
\def\Mor{\mathop{\mathrm{Mor}}\nolimits}
\def\Ob{\mathop{\mathrm{Ob}}\nolimits}
\def\Sh{\mathop{\mathit{Sh}}\nolimits}
\def\NL{\mathop{N\!L}\nolimits}
\def\CH{\mathop{\mathrm{CH}}\nolimits}
\def\proetale{{pro\text{-}\acute{e}tale}}
\def\etale{{\acute{e}tale}}
\def\QCoh{\mathit{QCoh}}
\def\Ker{\mathop{\mathrm{Ker}}}
\def\Im{\mathop{\mathrm{Im}}}
\def\Coker{\mathop{\mathrm{Coker}}}
\def\Coim{\mathop{\mathrm{Coim}}}

% Boxtimes
%
\DeclareMathSymbol{\boxtimes}{\mathbin}{AMSa}{"02}

%
% Macros for moduli stacks/spaces
%
\def\QCohstack{\mathcal{QC}\!\mathit{oh}}
\def\Cohstack{\mathcal{C}\!\mathit{oh}}
\def\Spacesstack{\mathcal{S}\!\mathit{paces}}
\def\Quotfunctor{\mathrm{Quot}}
\def\Hilbfunctor{\mathrm{Hilb}}
\def\Curvesstack{\mathcal{C}\!\mathit{urves}}
\def\Polarizedstack{\mathcal{P}\!\mathit{olarized}}
\def\Complexesstack{\mathcal{C}\!\mathit{omplexes}}
% \Pic is the operator that assigns to X its picard group, usage \Pic(X)
% \Picardstack_{X/B} denotes the Picard stack of X over B
% \Picardfunctor_{X/B} denotes the Picard functor of X over B
\def\Pic{\mathop{\mathrm{Pic}}\nolimits}
\def\Picardstack{\mathcal{P}\!\mathit{ic}}
\def\Picardfunctor{\mathrm{Pic}}
\def\Deformationcategory{\mathcal{D}\!\mathit{ef}}

% --- appended 2026-10-12 (Chief): the project own definitions of the reference/slogan/history
% display environments, copied mechanically from the shared Stacks preamble lines 8-12
% (\usepackage{verbatim} plus the three \newenvironment lines, which the chapter preamble keeps
% verbatim from the project preamble). Nothing added, removed or redefined: these environments
% carry no mathematical content and the project itself typesets them as comments. Without these
% three lines a counted statement whose own text uses them fails with Environment slogan undefined.
\usepackage{verbatim}
\newenvironment{reference}{\comment}{\endcomment}
\newenvironment{slogan}{\comment}{\endcomment}
\newenvironment{history}{\comment}{\endcomment}
\usepackage{cleveref}
\crefname{originallemma}{original lemma}{original lemmas}
\Crefname{originallemma}{Original lemma}{Original lemmas}
\makeatletter
\newlabel{more-algebra-lemma-map-P-ring-to-completion-P}{{0BK9}{source}{}{sourcepointer-1}{}}
\newlabel{more-algebra-lemma-map-P-ring-to-completion-P@cref}{{[originallemma][0][]0BK9}{[1][0][]source}{}{}{}}
\makeatother
\crefname{originallemma}{original lemma}{original lemmas}
\Crefname{originallemma}{Original lemma}{Original lemmas}
\makeatletter
\newlabel{algebra-lemma-completion-flat}{{00MB}{source}{}{sourcepointer-2}{}}
\newlabel{algebra-lemma-completion-flat@cref}{{[originallemma][0][]00MB}{[1][0][]source}{}{}{}}
\makeatother
\crefname{originallemma}{original lemma}{original lemmas}
\Crefname{originallemma}{Original lemma}{Original lemmas}
\makeatletter
\newlabel{algebra-lemma-completion-Noetherian-Noetherian}{{0316}{source}{}{sourcepointer-3}{}}
\newlabel{algebra-lemma-completion-Noetherian-Noetherian@cref}{{[originallemma][0][]0316}{[1][0][]source}{}{}{}}
\makeatother
\crefname{originallemma}{original lemma}{original lemmas}
\Crefname{originallemma}{Original lemma}{Original lemmas}
\makeatletter
\newlabel{lemma-P-local}{{0BK8}{source}{}{sourcepointer-4}{}}
\newlabel{lemma-P-local@cref}{{[originallemma][0][]0BK8}{[1][0][]source}{}{}{}}
\makeatother
\crefname{originallemma}{original lemma}{original lemmas}
\Crefname{originallemma}{Original lemma}{Original lemmas}
\makeatletter
\newlabel{algebra-lemma-radical-completion}{{05GI}{source}{}{sourcepointer-5}{}}
\newlabel{algebra-lemma-radical-completion@cref}{{[originallemma][0][]05GI}{[1][0][]source}{}{}{}}
\makeatother
\crefname{originallemma}{original lemma}{original lemmas}
\Crefname{originallemma}{Original lemma}{Original lemmas}
\makeatletter
\newlabel{lemma-flat-unramified}{{0AGX}{source}{}{sourcepointer-6}{}}
\newlabel{lemma-flat-unramified@cref}{{[originallemma][0][]0AGX}{[1][0][]source}{}{}{}}
\makeatother
\crefname{originallemma}{original lemma}{original lemmas}
\Crefname{originallemma}{Original lemma}{Original lemmas}
\makeatletter
\newlabel{algebra-lemma-completion-faithfully-flat}{{00MC}{source}{}{sourcepointer-7}{}}
\newlabel{algebra-lemma-completion-faithfully-flat@cref}{{[originallemma][0][]00MC}{[1][0][]source}{}{}{}}
\makeatother
\title{The fibres of the completion map of a P-ring have property P}
\author{The Stacks Project Authors}
\date{stacks-project:more-algebra.tex:local-snapshot-20261012; GNU FDL 1.2 or later}
\begin{document}
\maketitle
\setcounter{subsection}{1}
\noindent\textit{Editorial note.} All delivered mathematical text is the author's own, copied from the pinned Stacks Project chapter file more-algebra.tex; no statement or proof step was rewritten, reordered or removed.\par\medskip
\noindent\textit{Editorial note.} The chapter file pinned-original-source.tex carries only \textbackslash{}input\{preamble\} on its first line, so no preamble text lives inside the chapter. The retained file sources/stacks-more-algebra/retained-preamble.tex is a mechanical, line-for-line copy of the project's single shared preamble as staged in the snapshot (Xy-pic setup plus theorem declarations and mathematical macros). Nothing was added, removed or redefined, and the excerpt serves only as preamble, never as a span coordinate.\par\medskip
\noindent\textit{Editorial note.} Both counted spans start at the \textbackslash{}begin line of their environment, so the statement and the author's proof are delivered whole; surrounding prose of the section is deliberately outside the counted spans and is carried as standard prerequisites instead.\par\medskip
\noindent\textit{Editorial note.} The citation fields (source\_id, source\_version, authors, source\_url, license.name, source\_title) are identical to those of every other card this lane has produced from this chapter, so the batch selector treats them as one source.\par\medskip
\noindent\textit{Allowed standard prerequisites (editorial).}\begin{itemize}
\item Flatness and Noetherianity of completions. Algebra, Lemma \textbackslash{}ref\{algebra-lemma-completion-flat\} and Algebra, Lemma \textbackslash{}ref\{algebra-lemma-completion-Noetherian-Noetherian\}, cited at proof lines 13055-13058. Let the author verify the local criterion for property P for the completion map.
\item The local criterion for property P: it is enough to check the third condition on a flat map whose source is Noetherian. Lemma \textbackslash{}ref\{lemma-P-local\}, cited at proof lines 13059-13060. Reduces the claim to the behaviour at a maximal ideal of the completion.
\item The ideal of definition is contained in every maximal ideal of the completion, and the quotient of the completion by it is the quotient of the ring. Algebra, Lemma \textbackslash{}ref\{algebra-lemma-radical-completion\}, cited at proof line 13063. Identifies the residue fields of the two local rings and shows that the prime below is maximal.
\item A flat local map of Noetherian local rings identifying residue fields and with equal maximal ideals induces an isomorphism on completions. Lemma \textbackslash{}ref\{lemma-flat-unramified\}, cited at proof lines 13073-13074. Identifies the completion of the localised completion with the completion of the localisation.
\item The completion map of a local ring is faithfully flat, and property P for fibres is preserved by the composition witnessed by a faithfully flat map. Algebra, Lemma \textbackslash{}ref\{algebra-lemma-completion-faithfully-flat\}, cited at proof lines 13077-13079, and condition (D), used at proof lines 13080-13083. Allow the permanence condition (D) to be applied to the two-step chain of ring maps.
\end{itemize}\medskip
\ifdefined\hypertarget\hypertarget{sourcepointer-1}{}\fi
\noindent\textit{Source reference.} Original lemma 0BK9: focused-seeds/AGP/sources/stacks-snapshot-20261011/tags.txt, source lines 14552--14552. The lemma's own canonical tag, quoted so that readers can resolve the label inside the delivered text; it points at the counted statement itself and carries no proof content.\par\medskip
\ifdefined\hypertarget\hypertarget{sourcepointer-2}{}\fi
\noindent\textit{Source reference.} Original lemma 00MB: focused-seeds/AGP/sources/stacks-snapshot-20261011/tags.txt, source lines 817--817. Navigation pointer to a statement the counted proof refers to but does not contain; nothing load-bearing is delegated to it.\par\medskip
\ifdefined\hypertarget\hypertarget{sourcepointer-3}{}\fi
\noindent\textit{Source reference.} Original lemma 0316: focused-seeds/AGP/sources/stacks-snapshot-20261011/tags.txt, source lines 3973--3973. Navigation pointer to a statement the counted proof refers to but does not contain; nothing load-bearing is delegated to it.\par\medskip
\ifdefined\hypertarget\hypertarget{sourcepointer-4}{}\fi
\noindent\textit{Source reference.} Original lemma 0BK8: focused-seeds/AGP/sources/stacks-snapshot-20261011/tags.txt, source lines 14551--14551. Navigation pointer to a statement the counted proof refers to but does not contain; nothing load-bearing is delegated to it.\par\medskip
\ifdefined\hypertarget\hypertarget{sourcepointer-5}{}\fi
\noindent\textit{Source reference.} Original lemma 05GI: focused-seeds/AGP/sources/stacks-snapshot-20261011/tags.txt, source lines 7062--7062. Navigation pointer to a statement the counted proof refers to but does not contain; nothing load-bearing is delegated to it.\par\medskip
\ifdefined\hypertarget\hypertarget{sourcepointer-6}{}\fi
\noindent\textit{Source reference.} Original lemma 0AGX: focused-seeds/AGP/sources/stacks-snapshot-20261011/tags.txt, source lines 13210--13210. Navigation pointer to a statement the counted proof refers to but does not contain; nothing load-bearing is delegated to it.\par\medskip
\ifdefined\hypertarget\hypertarget{sourcepointer-7}{}\fi
\noindent\textit{Source reference.} Original lemma 00MC: focused-seeds/AGP/sources/stacks-snapshot-20261011/tags.txt, source lines 818--818. Navigation pointer to a statement the counted proof refers to but does not contain; nothing load-bearing is delegated to it.\par\medskip
\begin{lemma}
\label{lemma-map-P-ring-to-completion-P}
Let $A$ be a $P$-ring where $P$ satisfies (B) and (D).
Let $I \subset A$ be an ideal and let $A^\wedge$ be the completion of $A$
with respect to $I$. Then the fibres of $A \to A^\wedge$ have $P$.
\end{lemma}

\begin{proof}
The ring map $A \to A^\wedge$ is flat by
Algebra, Lemma \ref{algebra-lemma-completion-flat}.
The ring $A^\wedge$ is Noetherian by
Algebra, Lemma \ref{algebra-lemma-completion-Noetherian-Noetherian}.
Thus it suffices to check the third condition of
Lemma \ref{lemma-P-local}.
Let $\mathfrak m' \subset A^\wedge$ be a maximal ideal lying over
$\mathfrak m \subset A$.
By Algebra, Lemma \ref{algebra-lemma-radical-completion}
we have $IA^\wedge \subset \mathfrak m'$.
Since $A^\wedge/IA^\wedge = A/I$ we see that
$I \subset \mathfrak m$, $\mathfrak m/I = \mathfrak m'/IA^\wedge$, and
$A/\mathfrak m = A^\wedge/\mathfrak m'$. Since $A^\wedge/\mathfrak m'$
is a field, we conclude that $\mathfrak m$ is a maximal ideal as well.
Then $A_\mathfrak m \to A^\wedge_{\mathfrak m'}$ is a flat local
ring homomorphism of Noetherian local rings
which identifies residue fields and such that
$\mathfrak m A^\wedge_{\mathfrak m'} = \mathfrak m'A^\wedge_{\mathfrak m'}$.
Thus it induces an isomorphism on complete local rings, see
Lemma \ref{lemma-flat-unramified}.
Let $(A_\mathfrak m)^\wedge$ be the completion of $A_\mathfrak m$
with respect to its maximal ideal.
The ring map
$$
(A^\wedge)_{\mathfrak m'} \to
((A^\wedge)_{\mathfrak m'})^\wedge = (A_\mathfrak m)^\wedge
$$
is faithfully flat (Algebra, Lemma
\ref{algebra-lemma-completion-faithfully-flat}). Thus we can apply
(D) to the ring maps
$$
A_\mathfrak m \to (A^\wedge)_{\mathfrak m'} \to (A_\mathfrak m)^\wedge
$$
to conclude because the fibres of
$A_\mathfrak m \to (A_\mathfrak m)^\wedge$
have $P$ as $A$ is a $P$-ring.
\end{proof}

\end{document}
